% Faithful formal-publication transcription of the main statements and necessary query-model definitions.
\section*{Author result and reduction}
We assume $n\ge3$ for simplicity of notation throughout the paper.
\begin{authorclaim}{Theorem2 (Maj lemma)}
$R(\mathrm{Maj}\circ g^n)=\Theta(n\cdot\overline R_{1/n}(g))$ for every partial function $g$.
\end{authorclaim}
Previously, Ben-David et al. [5] proved Theorem2 in the special case $g=\mathrm{GapOr}$. Here $\mathrm{GapOr}=\mathrm{GapOr}_m$ is the $m$-bit partial function defined by $\mathrm{GapOr}(x)=\mathrm{Or}(x)$ on inputs of Hamming weight $|x|\in\{0,m/2\}$ and is undefined otherwise. This is a particularly clean example of a function whose query complexity behaves as (assuming $m\ge\log(1/\epsilon)$)
\[\overline R_\epsilon(\mathrm{GapOr})=\Theta(\log(1/\epsilon)).\]
We prove Theorem2 by a direct reduction to this previous result! Our more general result says, informally, that error reduction is necessary for any composed function $f\circ g^n$ if it is necessary in the special case $g=\mathrm{GapOr}$. Our key conceptual insight is to formulate a sense in which every $g$ can be ``simulated'' by GapOr. There is, however, a slight technical caveat. For the reduction to work, we need to assume that the lower bound for $f\circ\mathrm{GapOr}^n$ holds not only against randomised decision trees but also against a more powerful model called $\epsilon$-approximate nonnegative degree $\deg^+_\epsilon$ (aka conical junta degree, partition bound), which we will recall in Section2.
\begin{authorclaim}{Theorem3 (Reduction to GapOr)}
If a function $f$ satisfies $\deg^+_\epsilon(f\circ h^n)\ge\Omega(n\log n)$ for some constant $\epsilon>0$ and for both $h\in\{\mathrm{GapOr}_{\log n},\neg\mathrm{GapOr}_{\log n}\}$, then
\[\forall g:\quad R(f\circ g^n)=\Omega(n\cdot\overline R_{1/n}(g)).\]
\end{authorclaim}
Theorem2 follows immediately by combining Theorem3 with [5, Theorem4], which proved the required nonnegative degree lower bound for $\mathrm{Maj}\circ\mathrm{GapOr}^n$ (we only note that their proof works equally well for $\neg\mathrm{GapOr}$ in place of GapOr). In fact, the nonnegative degree lower bound holds more generally for any $(2n+1)$-bit outer function that agrees with Maj on inputs of weight $n$ and $n+1$. For example, Xor is such a function, and hence the Xor lemma of Brody et al. [8] can be recovered using Theorem3. However, the original proof in [8] is much simpler than ours, and moreover, the result of [8] actually characterises $\overline R_\epsilon(\mathrm{Xor}\circ g^n)$ for all $\epsilon>0$ while we focus on the bounded-error case $\epsilon=1/3$.
Our goal for the rest of the paper is to prove Theorem3.
For randomised query complexity, a complete answer [to the direct sum problem] was recently obtained by Blais and Brody [7] (improving on [17,6]). Their result states
\[\forall g:\quad R(g^n)=\Theta(n\cdot\overline R_{1/n}(g)).\tag{Direct sum}\]
\section*{2 Query complexity basics}
We study partial boolean functions $f:\{0,1\}^n\to\{0,1,*\}$. The domain of the function is $\operatorname{dom}(f):=f^{-1}(\{0,1\})$ and the inputs $f^{-1}(*)$ are undefined. We say $f$ is total if $\operatorname{dom}(f)=\{0,1\}^n$. For partial functions $f$ and $g$, their composition $f\circ g^n$ is defined by $(f\circ g^n)(x_1,\ldots,x_n):=f(g(x_1),\ldots,g(x_n))$ if $x_i\in\operatorname{dom}(g)$ for all $i\in[n]$; otherwise $(f\circ g^n)(x_1,\ldots,x_n):=*$. Standard references for boolean function complexity are [9,18].
\paragraph{Decision trees.} A (deterministic) decision tree $t$ is an algorithm for computing a boolean function on an unknown input $x\in\{0,1\}^n$. The algorithm repeatedly queries the input variables $x_i\in\{0,1\}$ in some order (which can depend on outcomes of queries made so far) until eventually producing an output $t(x)$. Such an algorithm can be represented as a binary tree, with internal nodes labelled with variables $x_i$, outgoing edges of the internal nodes labelled with query outcomes ($x_i=0$ and $x_i=1$), and leaves labelled with output values. Each input $x$ determines a unique root-to-leaf path, obtained by following the query outcomes consistent with $x$. The most important cost measure of $t$ is its depth, denoted $\operatorname{depth}(t)$, which is the longest root-to-leaf path in the tree and equals $\max_xq(t,x)$ where $q(t,x)$ denotes the number of queries made by $t$ on input $x$.
A randomised decision tree $T$ is a distribution over deterministic decision trees $t\sim T$. We say $T$ computes $f:\{0,1\}^n\to\{0,1,*\}$ with error $\epsilon$ if for every $x\in\operatorname{dom}(f)$ we have $\Pr_{t\sim T}[t(x)=f(x)]\ge1-\epsilon$. There are two cost measures for $T$: the (worst-case) depth is the maximum depth of any decision tree in the support of $T$; the expected depth is $\max_x\mathbb E_{t\sim T}[q(t,x)]$. The $\epsilon$-error query complexity of $f$, denoted $R_\epsilon(f)$, is the least depth of a randomised decision tree that computes $f$ with error $\epsilon$. The $\epsilon$-error expected query complexity, denoted $\overline R_\epsilon(f)$, is defined analogously.
\paragraph{Error reduction.} It is well known that the error probability of an algorithm (computing a boolean-valued function) can be reduced from any constant $1/2-\delta$, where $\delta>0$, to any other constant $\epsilon>0$ by repeating the algorithm constantly many times (in fact, $O(\log(1/\epsilon)/\delta^2)$ many) and outputting the majority answer. Hence we often set $\epsilon:=1/3$ and omit $\epsilon$ from notation. In this bounded-error regime, we have $\overline R(f)\le R(f)\le O(\overline R(f))$ where the second inequality follows by truncating executions that query many more bits than the expectation. For vanishing $\epsilon=o(1)$ (as $n\to\infty$), it is possible that $\overline R_\epsilon(f)\le o(R_\epsilon(f))$. For example, consider the partial $2n$-bit function $f$ where the task is to distinguish inputs of the form $x0^n$ from inputs of the form $0^nx$ with the promise that $|x|=n/2$. We have $\overline R_{1/n}(f)=O(1)$ while $R_{1/n}(f)=\Theta(\log n)$. In this small-error regime, the following fine-grained error reduction calculation will be useful.
\begin{authorclaim}{Claim7}
$\overline R_{\epsilon^k}(f)\le4k\cdot\overline R_\epsilon(f)$ for every $k\ge1$ and $\epsilon\le1/16$.
\end{authorclaim}
\begin{proof}
Suppose $T$ computes $f$ with error $\epsilon$ and consider the algorithm $T'$ that runs $T$ $4k-1$ times and outputs the majority answer. Then $T'$ errs iff at least $2k$ of the runs err. This happens with probability at most $\sum_{i=2k}^{4k-1}\binom{4k-1}i\epsilon^i(1-\epsilon)^{4k-1-i}\le2^{4k}\epsilon^{2k}\le\epsilon^k$.
\end{proof}
\paragraph{Leaf indicators.} Let $t$ be a decision tree with $n$-bit inputs. We denote by $L(t)$ the set of its leaves and by $\ell_x^t\in L(t)$ the unique leaf reached on input $x$. We often identify a leaf $\ell\in L(t)$ with its associated leaf indicator function $\ell:\{0,1\}^n\to\{0,1\}$ defined by $\ell(x):=1$ iff input $x$ reaches leaf $\ell$. Thus each $\ell$ is simply a conjunction of at most $\operatorname{depth}(t)$ literals ($x_i$ or $\overline x_i$) determined by the unique root-to-$\ell$ path in $t$. If $t$ outputs boolean values, we let $A(t)\subseteq L(t)$ denote the set of accepting leaves, that is, those that output 1. Since the leaf indicators have pairwise disjoint supports, we can write the function computed by $t$ as
\begin{equation}\tag{1}t(x)=\sum_{\ell\in A(t)}\ell(x).\end{equation}
\paragraph{Nonnegative degree.} Let $p:\{0,1\}^n\to\mathbb R_{\ge0}$ be a nonnegative function. We say $p$ is a nonnegative $d$-junta if it depends on at most $d$ of its variables. For example, if $t$ is a depth-$d$ decision tree, then each $\ell\in L(t)$ is a nonnegative $d$-junta. More generally, we say that $p$ is a conical junta of degree $d$ if it can be written as a conical combination of nonnegative $d$-juntas, that is, $p(x)=\sum_i a_iq_i(x)$ where $a_i\ge0$ are nonnegative scalars and the $q_i$ are nonnegative $d$-juntas. For example, the function computed by $t$ is a degree-$d$ conical junta, as given by the expression(1). The nonnegative degree of $p$, denoted $\deg^+(p)$, is the least $d$ such that $p$ is a degree-$d$ conical junta.
Let $f:\{0,1\}^n\to\{0,1,*\}$ be a partial function. We say that $p$ $\epsilon$-approximates $f$ if $p(x)\in f(x)\pm\epsilon$ for every $x\in\operatorname{dom}(f)$. The $\epsilon$-approximate nonnegative degree of $f$, denoted $\deg^+_\epsilon(f)$, is the least degree of a conical junta that $\epsilon$-approximates $f$. For example, if $T$ is a depth-$d$ randomised $\epsilon$-error decision tree for $f$, then there exists a degree-$d$ conical junta $p_T$ that $\epsilon$-approximates $f$, namely,
\[p_T(x):=\mathbb E_{t\sim T}[t(x)]\in f(x)\pm\epsilon.\]
This shows that $\deg^+_\epsilon(f)\le R_\epsilon(f)$. The gap betweeen $\deg^+_{1/3}(f)$ and $R(f)$ can be huge for partial functions. For example, consider the $n$-bit UniqueOr defined by $\mathrm{UniqueOr}(x)=\mathrm{Or}(x)$ for inputs of weight $|x|\in\{0,1\}$ and undefined othwerwise. Then $\deg^+(\mathrm{UniqueOr})=1$ (computed by $\sum_ix_i$) while $R(\mathrm{UniqueOr})=\Theta(n)$. For total functions, the gap is at most polynomial [9].
Nonnegative degree has been studied under many names: (one-sided) partition bound [16], WAPP query complexity [14,5], and query complexity ``in expectation'' [19].
\section*{Necessary Leaf and Multileaf Lemma definitions (Section3.1)}
For a leaf $\ell\in L(t)$ and any input distribution $\mu$ we write for short $\ell(\mu):=\mathbb E_{x\sim\mu}[\ell(x)]=\Pr_{x\sim\mu}[\ell(x)=1]$.
Fix a partial function $g:\{0,1\}^m\to\{0,1,*\}$. Let $\mu=\frac12(\mu^0+\mu^1)$ be a balanced distribution where $\mu^i$ is supported on $g^{-1}(i)$. For a leaf $\ell$ over $m$ bits, a ``hard side'' $b\in\{0,1\}$, and an error parameter $\epsilon\ge0$, we define
\[\ell\text{ is }(\epsilon,\mu,b)\text{-noisy}\quad\Longleftrightarrow\quad\ell(\mu^{1-b})\ge\epsilon\cdot\ell(\mu^b).\]
For simplicity of notation, for small quantities $a,b\in[0,1]$, we write $a\ll b$ (resp. $a\lll b$) to mean $a\le cb$ (resp. $a^c\le b$) for a sufficiently small constant $c>0$.
\begin{authorclaim}{Leaf Lemma}
For every partial $g$ and $0<\epsilon\lll\delta\ll1$, there exists a distribution $\mu=\frac12(\mu^0+\mu^1)$ over $\operatorname{dom}(g)$ and a hard side $b\in\{0,1\}$ such that for every deterministic tree $t$ and $i\in\{0,1\}$:
\[\frac{\mathbb E_{x\sim\mu^i}[q(t,x)]}{\overline R_\epsilon(g)}\lll\delta\quad\Longrightarrow\quad\Pr_{x\sim\mu^i}[\ell_x^t\text{ is }(\epsilon,\mu,b)\text{-noisy}]\ge1-\delta.\]
\end{authorclaim}
In order to use Leaf Lemma in the context of composed functions, we generalise it to the direct sum setting where the inputs come from $\operatorname{dom}(g^n):=\operatorname{dom}(g)^n$. Let $\ell$ be a leaf over $nm$ bits and write $\ell(x)=\prod_{i\in[n]}\ell_i(x_i)$ where $x_i\in\{0,1\}^m$ and each $\ell_i$ is over $m$ bits. We define
\[\ell\text{ is }(\delta,\epsilon,\mu,b)\text{-noisy}\quad\Longleftrightarrow\quad\ell_i\text{ is }(\epsilon,\mu,b)\text{-noisy for at least }(1-\delta)n\text{ many }i\in[n].\]
Our generalised lemma says that we will typically reach a noisy leaf if we run a shallow decision tree on a random input from the product distribution $\mu^y:=\mu^{y_1}\times\cdots\times\mu^{y_n}$ where $y\in\{0,1\}^n$.
\begin{authorclaim}{Multileaf Lemma}
For every partial $g$ and $0<\epsilon\lll\delta\ll1$, there exists a distribution $\mu=\frac12(\mu^0+\mu^1)$ over $\operatorname{dom}(g)$ and a hard side $b\in\{0,1\}$ such that for every deterministic tree $t$ taking inputs from $\operatorname{dom}(g^n)$ and having $\operatorname{depth}(t)/(n\overline R_\epsilon(g))\lll\delta$,
\[\forall y\in\{0,1\}^n:\quad\Pr_{x\sim\mu^y}[\ell_x^t\text{ is }(\delta,\epsilon,\mu,b)\text{-noisy}]\ge1-\delta.\]
\end{authorclaim}
